% GENERATED FILE. Edit canonical lessons, then run npm run generate:books.
% canonical-source-hash: fnv1a64:a381d276f5da3d26
% canonical-lessons: ML-C478-vinodasancari, ML-C478-valikatti, ML-C478-laksyasthanam, ML-C478-yatraparipati, ML-C478-currikkarannuka

\chapter{Sightseeing}
\label{ch:ml-sightseeing}
\hlchaptermodality{478}

\begin{chapteropening}
\textbf{By the end of this chapter:} \emph{I can say a tourist, a guide, a destination, an itinerary and to roam around.}

Complete the last lesson of chapter 478: I can say a tourist, a guide, a destination, an itinerary and to roam around.
\end{chapteropening}

\section[\texorpdfstring{\ml{വിനോദസഞ്ചാരി} (vinōdasañcāri)}{വിനോദസഞ്ചാരി (vinōdasañcāri)}]{\ml{വിനോദസഞ്ചാരി} (vinōdasañcāri) — a tourist}

\label{lesson:ML-C478-vinodasancari}



\begin{quote}
Before the new one: say the Malayalam for to swim across, then the Malayalam for to dash forward.
\end{quote}

\subsection*{You'll want to know: \ml{വിനോദസഞ്ചാരി}}
\textbf{\ml{വിനോദസഞ്ചാരി}} — \emph{vinōdasañcāri} — \textquotedblleft{}a tourist\textquotedblright{}.

\begin{quote}
\ml{മഴക്കാലത്തും} \ml{ധാരാളം} \ml{വിനോദസഞ്ചാരികൾ} \ml{ഇവിടെ} \ml{വരുന്നു}.
\end{quote}

\begin{quote}
\emph{maḻakkālattuṁ dhārāḷaṁ vinōdasañcārikaḷ iviṭe varunnu.}
\end{quote}

\begin{quote}
Many tourists come here even in the rainy season.
\end{quote}

Plural: \ml{വിനോദസഞ്ചാരികൾ}, \emph{vinōdasañcārikaḷ}.

Literally a pleasure traveller: \ml{വിനോദം}, entertainment, plus \ml{സഞ്ചാരി}, a traveller.

\begin{cousinweb}
\textbf{\ml{വിനോദസഞ്ചാരി}} (\emph{vinōdasañcāri}) is a Malayalam compound built from two Sanskrit words, \emph{vinoda}, \textquotedblleft{}pastime, diversion\textquotedblright{}, and \emph{saṁcārin}, \textquotedblleft{}one who moves about\textquotedblright{} (Monier-Williams).
\end{cousinweb}

\begin{grammarlens}[title={What you've built}]
One more word for getting around.
\end{grammarlens}

\subsection*{Your turn}
\begin{itemize}
\raggedright
  \item \emph{Say it:} \emph{vinōdasañcāri}
  \item \emph{Say it:} \emph{vinōdasañcāri}, once more
  \item \emph{Say it:} \emph{maḻakkālattuṁ dhārāḷaṁ vinōdasañcārikaḷ iviṭe varunnu.} — Many tourists come here even in the rainy season.
  \item \emph{Recall:} say the Malayalam for to swim across, then the Malayalam for to dash forward, then say \emph{vinōdasañcāri} again
\end{itemize}

\begin{tcolorbox}[breakable,colback=teal!4,colframe=teal!35!black,title={Before you move on}]
What does \ml{വിനോദസഞ്ചാരി} mean? (\textquotedblleft{}A tourist\textquotedblright{}.) Say it once more.
\end{tcolorbox}
\section[\texorpdfstring{\ml{വഴികാട്ടി} (vaḻikāṭṭi)}{വഴികാട്ടി (vaḻikāṭṭi)}]{\ml{വഴികാട്ടി} (vaḻikāṭṭi) — a guide}

\label{lesson:ML-C478-valikatti}



\begin{quote}
Before the new one: say the Malayalam for to dash forward, then the Malayalam for a tourist.
\end{quote}

\subsection*{You'll want to know: \ml{വഴികാട്ടി}}
\textbf{\ml{വഴികാട്ടി}} — \emph{vaḻikāṭṭi} — \textquotedblleft{}a guide\textquotedblright{}.

\begin{quote}
\ml{ഒരു} \ml{വഴികാട്ടി} \ml{ഞങ്ങളെ} \ml{കാട്ടിലൂടെ} \ml{കൊണ്ടുപോയി}.
\end{quote}

\begin{quote}
\emph{oru vaḻikāṭṭi ñaṅṅaḷe kāṭṭilūṭe koṇṭupōyi.}
\end{quote}

\begin{quote}
A guide took us through the forest.
\end{quote}

Literally one who shows the way. A guidebook can also be a \ml{വഴികാട്ടി}.

\begin{cousinweb}
\textbf{\ml{വഴികാട്ടി}} (\emph{vaḻikāṭṭi}) is a native Dravidian compound of \ml{വഴി}, a way (DEDR 5297), and \ml{കാട്ടുക}, to show (DEDR 1443), with the ending -\ml{ഇ} for the one who does it.
\end{cousinweb}

\begin{grammarlens}[title={What you've built}]
One more word for getting around.
\end{grammarlens}

\subsection*{Your turn}
\begin{itemize}
\raggedright
  \item \emph{Say it:} \emph{vaḻikāṭṭi}
  \item \emph{Say it:} \emph{vaḻikāṭṭi}, once more
  \item \emph{Say it:} \emph{oru vaḻikāṭṭi ñaṅṅaḷe kāṭṭilūṭe koṇṭupōyi.} — A guide took us through the forest.
  \item \emph{Recall:} say the Malayalam for to dash forward, then the Malayalam for a tourist, then say \emph{vaḻikāṭṭi} again
\end{itemize}

\begin{tcolorbox}[breakable,colback=teal!4,colframe=teal!35!black,title={Before you move on}]
What does \ml{വഴികാട്ടി} mean? (\textquotedblleft{}A guide\textquotedblright{}.) Say it once more.
\end{tcolorbox}
\section[\texorpdfstring{\ml{ലക്ഷ്യസ്ഥാനം} (lakṣyasthānaṁ)}{ലക്ഷ്യസ്ഥാനം (lakṣyasthānaṁ)}]{\ml{ലക്ഷ്യസ്ഥാനം} (lakṣyasthānaṁ) — a destination}

\label{lesson:ML-C478-laksyasthanam}



\begin{quote}
Before the new one: say the Malayalam for a tourist, then the Malayalam for a guide.
\end{quote}

\subsection*{You'll want to know: \ml{ലക്ഷ്യസ്ഥാനം}}
\textbf{\ml{ലക്ഷ്യസ്ഥാനം}} — \emph{lakṣyasthānaṁ} — \textquotedblleft{}a destination\textquotedblright{}.

\begin{quote}
\ml{രാത്രി} \ml{എട്ട്} \ml{മണിക്ക്} \ml{ഞങ്ങൾ} \ml{ലക്ഷ്യസ്ഥാനത്ത്} \ml{എത്തി}.
\end{quote}

\begin{quote}
\emph{rātri eṭṭŭ maṇikkŭ ñaṅṅaḷ lakṣyasthānattŭ etti.}
\end{quote}

\begin{quote}
We reached our destination at eight at night.
\end{quote}

\ml{ലക്ഷ്യം}, an aim, plus \ml{സ്ഥാനം}, a place: the place you aim at.

\begin{cousinweb}
\textbf{\ml{ലക്ഷ്യസ്ഥാനം}} (\emph{lakṣyasthānaṁ}) is a Malayalam compound built from two words borrowed from Sanskrit, \emph{lakṣya}, \textquotedblleft{}a mark, an aim\textquotedblright{}, and \emph{sthāna}, \textquotedblleft{}a place\textquotedblright{} (Monier-Williams).
\end{cousinweb}

\begin{grammarlens}[title={What you've built}]
One more word for getting around.
\end{grammarlens}

\subsection*{Your turn}
\begin{itemize}
\raggedright
  \item \emph{Say it:} \emph{lakṣyasthānaṁ}
  \item \emph{Say it:} \emph{lakṣyasthānaṁ}, once more
  \item \emph{Say it:} \emph{rātri eṭṭŭ maṇikkŭ ñaṅṅaḷ lakṣyasthānattŭ etti.} — We reached our destination at eight at night.
  \item \emph{Recall:} say the Malayalam for a tourist, then the Malayalam for a guide, then say \emph{lakṣyasthānaṁ} again
\end{itemize}

\begin{tcolorbox}[breakable,colback=teal!4,colframe=teal!35!black,title={Before you move on}]
What does \ml{ലക്ഷ്യസ്ഥാനം} mean? (\textquotedblleft{}A destination\textquotedblright{}.) Say it once more.
\end{tcolorbox}
\section[\texorpdfstring{\ml{യാത്രാപരിപാടി} (yātrāparipāṭi)}{യാത്രാപരിപാടി (yātrāparipāṭi)}]{\ml{യാത്രാപരിപാടി} (yātrāparipāṭi) — an itinerary; a travel plan}

\label{lesson:ML-C478-yatraparipati}



\begin{quote}
Before the new one: say the Malayalam for a guide, then the Malayalam for a destination.
\end{quote}

\subsection*{You'll want to know: \ml{യാത്രാപരിപാടി}}
\textbf{\ml{യാത്രാപരിപാടി}} — \emph{yātrāparipāṭi} — \textquotedblleft{}an itinerary; a travel plan\textquotedblright{}.

\begin{quote}
\ml{മഴ} \ml{കാരണം} \ml{ഞങ്ങൾ} \ml{യാത്രാപരിപാടി} \ml{മാറ്റി}.
\end{quote}

\begin{quote}
\emph{maḻa kāraṇaṁ ñaṅṅaḷ yātrāparipāṭi māṟṟi.}
\end{quote}

\begin{quote}
Because of the rain we changed our travel plans.
\end{quote}

Built from \ml{യാത്ര}, a journey, and \ml{പരിപാടി}, a programme.

\begin{cousinweb}
\textbf{\ml{യാത്രാപരിപാടി}} (\emph{yātrāparipāṭi}) is a Malayalam compound built from \ml{യാത്ര}, from Sanskrit \emph{yātrā} (CDIAL 10456), and \ml{പരിപാടി}, a programme, from Sanskrit \emph{paripāṭī}, \textquotedblleft{}order, method\textquotedblright{} (Monier-Williams).
\end{cousinweb}

\begin{grammarlens}[title={What you've built}]
One more word for getting around.
\end{grammarlens}

\subsection*{Your turn}
\begin{itemize}
\raggedright
  \item \emph{Say it:} \emph{yātrāparipāṭi}
  \item \emph{Say it:} \emph{yātrāparipāṭi}, once more
  \item \emph{Say it:} \emph{maḻa kāraṇaṁ ñaṅṅaḷ yātrāparipāṭi māṟṟi.} — Because of the rain we changed our travel plans.
  \item \emph{Recall:} say the Malayalam for a guide, then the Malayalam for a destination, then say \emph{yātrāparipāṭi} again
\end{itemize}

\begin{tcolorbox}[breakable,colback=teal!4,colframe=teal!35!black,title={Before you move on}]
What does \ml{യാത്രാപരിപാടി} mean? (\textquotedblleft{}An itinerary; a travel plan\textquotedblright{}.) Say it once more.
\end{tcolorbox}
\section[\texorpdfstring{\ml{ചുറ്റിക്കറങ്ങുക} (cuṟṟikkaṟaṅṅuka)}{ചുറ്റിക്കറങ്ങുക (cuṟṟikkaṟaṅṅuka)}]{\ml{ചുറ്റിക്കറങ്ങുക} (cuṟṟikkaṟaṅṅuka) — to roam around; to go sightseeing}

\label{lesson:ML-C478-currikkarannuka}



\begin{quote}
Before the new one: say the Malayalam for a destination, then the Malayalam for an itinerary.
\end{quote}

\subsection*{You'll want to know: \ml{ചുറ്റിക്കറങ്ങുക}}
\textbf{\ml{ചുറ്റിക്കറങ്ങുക}} — \emph{cuṟṟikkaṟaṅṅuka} — \textquotedblleft{}to roam around; to go sightseeing\textquotedblright{}.

\begin{quote}
\ml{അവധിക്ക്} \ml{ഞങ്ങൾ} \ml{നഗരം} \ml{മുഴുവൻ} \ml{ചുറ്റിക്കറങ്ങി}.
\end{quote}

\begin{quote}
\emph{avadhikkŭ ñaṅṅaḷ nagaraṁ muḻuvan cuṟṟikkaṟaṅṅi.}
\end{quote}

\begin{quote}
In the holidays we went all round the city.
\end{quote}

Past: \ml{ചുറ്റിക്കറങ്ങി}, \emph{cuṟṟikkaṟaṅṅi}.

\ml{ചുറ്റി}, going round, plus \ml{കറങ്ങുക}, to circle. What is toured takes no ending.

\begin{cousinweb}
\textbf{\ml{ചുറ്റിക്കറങ്ങുക}} (\emph{cuṟṟikkaṟaṅṅuka}) is a native Dravidian compound of \ml{ചുറ്റുക}, to go round (DEDR 2715), and \ml{കറങ്ങുക}, to circle (DEDR 1387).
\end{cousinweb}

\begin{grammarlens}[title={What you've built}]
One more word for getting around.
\end{grammarlens}

\subsection*{Your turn}
\begin{itemize}
\raggedright
  \item \emph{Say it:} \emph{cuṟṟikkaṟaṅṅuka}
  \item \emph{Say it:} \emph{cuṟṟikkaṟaṅṅuka}, once more
  \item \emph{Say it:} \emph{avadhikkŭ ñaṅṅaḷ nagaraṁ muḻuvan cuṟṟikkaṟaṅṅi.} — In the holidays we went all round the city.
  \item \emph{Recall:} say the Malayalam for a destination, then the Malayalam for an itinerary, then say \emph{cuṟṟikkaṟaṅṅuka} again
\end{itemize}

\begin{tcolorbox}[breakable,colback=teal!4,colframe=teal!35!black,title={Before you move on}]
What does \ml{ചുറ്റിക്കറങ്ങുക} mean? (\textquotedblleft{}To roam around; to go sightseeing\textquotedblright{}.) Say it once more.
\end{tcolorbox}
